\chapter{Participant results}
\label{ch:participant}

\todo[inline]{WRITE: this chapter is a SKELETON. Every figure and table below
is a placeholder with the analysis script that produces it named in the
caption, so the chapter can be filled the moment data exists. No participant
has been run: the direct-teleoperation baseline depends on master channels
that are currently incoherent, and \cref{ch:discussion} lists what must be
resolved first.}

\section{Participants}

\begin{table}[htbp]
  \centering
  \caption[Demographics]{Demographics, by role. \textbf{Produced by}
  \repo{srl\_experiments/analyse\_demographics.py}. Reported per ROLE because
  operator and wearer are different people with different exposure: only the
  wearer carries \SI{17}{\kilo\gram}. No identifying field is stored;
  \repo{write\_manifest()} raises on name, email, date of birth, address or
  phone.}
  \label{tab:demographics}
  \begin{tabular}{@{}lccc@{}}
    \toprule
    & Operators & Wearers & Dyads \\
    \midrule
    n & -- & -- & -- \\
    age (median, IQR) & -- & -- & -- \\
    prior teleoperation experience & -- & -- & -- \\
    handedness & -- & -- & -- \\
    \bottomrule
  \end{tabular}
\end{table}

\section{Task performance}

\begin{figure}[htbp]
  \centering
  \imgplaceholder{55mm}{Completion time per task, by mode. Five panels (T1
  positioning, T2 pick-and-place, T3 rigid carry, T4 compliant carry, T5 dual
  pursuit), box plot per mode within each. \textbf{T2 CARRIES NO MODE
  COMPARISON} and must be drawn as a single condition: teleoperated grasping
  is not achievable on this platform, so it has no baseline.}
  \caption[Completion time by task and mode]{Completion time. \textbf{Produced
  by} \repo{analyse\_autonomy\_level.py}. Omnibus plus Holm-corrected post-hoc
  from \repo{srl\_experiments/posthoc.py}; report raw and adjusted $p$ side by
  side.}
  \label{fig:completion}
\end{figure}

\begin{figure}[htbp]
  \centering
  \imgplaceholder{50mm}{Coordination error traces for the coupled tasks: tray
  tilt (T3) and gripper separation error (T4) against time, one trace per
  condition, with the failure threshold drawn as a horizontal line
  (\SI{6.8}{\degree} tilt, \SI{534}{\milli\metre} separation).}
  \caption[Coordination error]{Coordination error over a trial.
  \textbf{Produced by} \repo{experiments/bimanual/coupled\_metrics.py}.}
  \label{fig:coordination}
\end{figure}

\begin{figure}[htbp]
  \centering
  \imgplaceholder{50mm}{Cross-arm interference. RMS tracking error of one arm
  against the OTHER arm's target speed, with the fitted linear model and
  $b_{\text{cross}}$ annotated. Report \texttt{phase\_lag\_s} beside it: a
  large $b_{\text{cross}}$ with a flat phase lag is attentional, both rising
  together is not.}
  \caption[Cross-arm interference]{Cross-arm interference in T5.
  \textbf{Produced by} \repo{analyse\_divided\_attention.py}. The coefficient
  is TOTAL speed sensitivity; amplitude and lag are not separable from RMS
  error alone.}
  \label{fig:interference}
\end{figure}

\begin{table}[htbp]
  \centering
  \caption[Grasp success]{Grasp success, T2 only, \textbf{modes 4 and above}.
  \textbf{Produced by} \repo{analyse\_autonomy\_level.py}. No teleoperation
  row exists and none should be added: see
  \repo{docs/research/09\_task2\_grasping\_finding.md}.}
  \label{tab:grasp}
  \begin{tabular}{@{}lccc@{}}
    \toprule
    Mode & attempts & successes & rate \\
    \midrule
    04 shared autonomy & -- & -- & -- \\
    05 supervised & -- & -- & -- \\
    06 full autonomy & -- & -- & -- \\
    \bottomrule
  \end{tabular}
\end{table}

\section{Workload and exertion}

\begin{figure}[htbp]
  \centering
  \imgplaceholder{48mm}{NASA-TLX, six subscales plus the weighted total, as
  grouped bars per mode. \textbf{Two panels, operator and wearer}, because the
  two load differently and pooling them would hide the effect the design is
  built to find.}
  \caption[NASA-TLX]{Workload~\cite{hart1988nasatlx}. \textbf{Produced by}
  \repo{analyse\_workload.py}.}
  \label{fig:tlx}
\end{figure}

\begin{figure}[htbp]
  \centering
  \imgplaceholder{42mm}{Borg CR10 against block index, one line per
  participant, operator and wearer separately. Fatigue is monotonic in time,
  so the SLOPE is the quantity of interest, not the endpoint.}
  \caption[Perceived exertion]{Borg CR10~\cite{borg1982perceived} across
  blocks. \textbf{Produced by} \repo{analyse\_workload.py}. Probed between
  blocks so fatigue is a covariate rather than an unmeasured confound.}
  \label{fig:borg}
\end{figure}

\section{Trust, embodiment and proprioceptive drift}

\begin{figure}[htbp]
  \centering
  \imgplaceholder{48mm}{Wearer trust by mode, with the SRL Proxemics finding
  drawn as a reference line. That work found higher autonomy LOWERED trust;
  this figure either replicates that or does not, and the reference line makes
  which visible at a glance.}
  \caption[Wearer trust]{Trust by autonomy level, wearer only.
  \textbf{Produced by} \repo{analyse\_trust.py}. Compare
  against~\cite{zhou2026proxemics}.}
  \label{fig:trust}
\end{figure}

\begin{figure}[htbp]
  \centering
  \imgplaceholder{48mm}{Embodiment, decomposed into ownership, agency and
  location rather than reported as one score, since the components move
  independently~\cite{longo2008embodiment}. Wearer and operator separately.}
  \caption[Embodiment]{Embodiment components. \textbf{Produced by}
  \repo{analyse\_embodiment.py}.}
  \label{fig:embodiment}
\end{figure}

\begin{figure}[htbp]
  \centering
  \imgplaceholder{42mm}{Proprioceptive drift, pre against post, one point per
  participant with the unity line drawn. Displacement from the line is the
  effect; the paired structure must be visible rather than summarised to two
  means.}
  \caption[Proprioceptive drift]{Proprioceptive drift, pre and post.
  \textbf{Produced by} \repo{analyse\_embodiment.py}.}
  \label{fig:drift}
\end{figure}

\todo[inline]{CONFIRM: whether the exchange rate between operator performance
and wearer trust can be estimated at $n=16$ dyads, or whether it is reported
descriptively with a confidence interval and no test. The design was powered
for $d=0.73$ on the within-subjects performance contrast, NOT for a
correlation across roles.}
